% % Abstract

\thispagestyle{empty}
\pdfbookmark[0]{Abstract}{Abstract} % Bookmark name visible in a PDF viewer

\begin{center}
%	\bigskip

    {\normalsize \href{http://www.sussex.ac.uk/}{\myUni} \\} % University name in capitals
    {\normalsize \myFaculty \\} % Faculty name
    {\normalsize \myDepartment \\} % Department name
    \bigskip\vspace*{.02\textheight}
    {\Large \textsc{Doctoral Thesis}}\par
    \bigskip
    
    {\rule{\linewidth}{1pt}\\%[0.4cm]
    \Large \myTitle \par} % Thesis title
    \rule{\linewidth}{1pt}\\[0.4cm]
    
    \bigskip
	{\normalsize by \myName \par} % Author name
    \bigskip\vspace*{.06\textheight}
\end{center}

    {\centering\Huge \textbf{Abstract} \par}
    \bigskip

\noindent
This thesis presents two research projects developing theoretical tools for the
description of QCD radiation in high-energy collisions, with applications to
jet physics at the LHC and lepton--hadron colliders.

The first project introduces an inclusive generalised-$k_t$ jet algorithm for
Deep Inelastic Scattering (DIS), defined in the Breit frame and implemented in
\texttt{fjcontrib}, a third-party extension of \textsc{FastJet}. The family of
algorithms is governed by the usual parameter $p$, which controls the
transverse-momentum dependence of the algorithm, together with a jet radius
parameter $R$. We investigate phenomenological applications related to the
identification of the jet associated with the struck quark, and assess their
sensitivity to non-perturbative effects such as hadronisation, alongside
comparisons with the recent Centauro algorithm.

The second extends the \textsc{ARES} method, a semi-numerical framework for
next-to-next-to-leading logarithmic (NNLL) resummation, to hadronic collisions,
considering events with a $Z$ boson recoiling against a high-transverse-momentum
jet. We find that the structure of the NNLL resummation is formally identical to
that for three-jet events in $e^+e^-$ annihilation, but with different
ingredients. In particular, initial-state radiation is incorporated within the
\textsc{ARES} framework for the first time, with the associated NNLL effects
involving the convolution of splitting functions with parton distribution
functions. The results are validated analytically against the known resummation
of one-jettiness.
 
